\chapter{Węzły i rozmaitości czterowymiarowe}
\label{ch:kt-cztery}

Chirurgia Dehna opisuje zamknięte trójrozmaitości jako brzegi
czterorozmaitości zbudowanych z uchwytów. Ta obserwacja jest
punktem wyjścia rachunku Kirby'ego, teorii form przecięcia,
problemu rozpoznawania egzotycznych struktur gładkich oraz
teorii kobordyzmu węzłów. Poniżej precyzujemy konstrukcje
od podstaw.

\section{Uchwyt dwuwymiarowy w wymiarze cztery}
\label{sec:kt-uchwyty}

\begin{definicja}[Uchwyt indeksu dwa]
\emph{Uchwytem indeksu dwa w wymiarze cztery} nazywamy
$h^2=D^2\times D^2$. Jego \emph{obszar przyczepienia}
to $\partial D^2\times D^2\cong S^1\times D^2$,
a \emph{rdzeń} to $D^2\times\{0\}$.
Dysk $\{0\}\times D^2$ nazywamy \emph{współrdzeniem}.
Rdzeń i współrdzeń przecinają się poprzecznie
w jednym punkcie; przy orientacji iloczynowej
znak tego przecięcia jest dodatni.
Jeżeli $W$ jest zwartą, orientowaną, gładką czterorozmaitością
z brzegiem, to dołączenie uchwytu $h^2$ do $W$ wymaga
wyboru gładkiego zanurzenia z ramowaniem
$\phi\colon S^1\times D^2\hookrightarrow\partial W$,
odwracającego orientację obszaru przyczepienia
jako części brzegu uchwytu względem orientacji
brzegu $W$. Dzięki temu orientacje sklejają się.
Wynikiem jest $W\cup_\phi(D^2\times D^2)$, po wygładzeniu
narożników.
\end{definicja}

Gdy $W=B^4$, krzywa $\phi(S^1\times\{0\})\subset\partial B^4=S^3$
jest węzłem $K$. Ramowanie decyduje o sposobie utożsamienia
otoczenia $K$ z $S^1\times D^2$. Tak więc \emph{węzeł wraz
z liczbą całkowitą} koduje konkretny sposób dołączenia
uchwytu indeksu dwa.

\begin{twierdzenie}[Brzeg po dołączeniu uchwytu]
\label{thm:kt-handle-boundary}
Niech $K\subset S^3=\partial B^4$ będzie zorientowanym węzłem,
a $n\in\mathbb Z$ jego całkowitym ramowaniem względem długości
preferowanej. Niech
$W_n(K):=B^4\cup_{(K,n)}(D^2\times D^2)$.
Wtedy $\partial W_n(K)$, z orientacją brzegową, jest
dyfeomorficzny z zachowaniem orientacji z rozmaitością $S^3_n(K)$
otrzymaną przez integralną chirurgię Dehna o współczynniku $n$.
\end{twierdzenie}

\begin{proof}
Brzeg iloczynu $D^2\times D^2$ składa się z dwóch
torusów pełnych $\partial D^2\times D^2$ i
$D^2\times\partial D^2$, spotykających się na torusie
$\partial D^2\times\partial D^2$. Pierwszy z nich
jest obszarem przyczepienia; jego wnętrze po sklejeniu
znika z brzegu. Drugi zastępuje więc otoczenie $N(K)$
w $S^3$. Dla ustalonego $z\in\partial D^2$ meridian
drugiego torusa ma postać $\partial D^2\times\{z\}$.
Jego obraz przez $\phi$ jest równoległą kopią $K$
określoną przez ramowanie, czyli ma klasę
$n[\mu]+[\lambda]$. Sklejenie na wspólnym torusie
odwraca orientacje brzegowe, ponieważ orientacja
całego dołączonego uchwytu zgadza się z orientacją
$B^4$ po sklejeniu. Po wygładzeniu narożników
otrzymujemy dokładnie gładkie wypełnienie Dehna
o nachyleniu $n/1$ (por. sekcję \ref{sec:kt-dehn}).
W szczególności dla splotu o rozłącznych otoczeniach
przyczepienia analogiczny opis daje jednoczesną
integralną chirurgię na wszystkich składowych.
\end{proof}

\begin{center}
\begin{tikzpicture}[>=Stealth,scale=.90,line cap=round]
  \node[draw=ktBlue,line width=1.5pt,rounded corners=8pt,
        minimum width=2.4cm,minimum height=1.32cm,
        fill=ktBlue!10,text=ktInk] (b) at (-3.0,0) {$B^4$};
  \node[draw=ktPurple,line width=1.5pt,rounded corners=8pt,
        minimum width=3.2cm,minimum height=1.32cm,
        fill=ktPurple!10,text=ktInk] (h) at (3.0,0) {$D^2\times D^2$};
  \draw[->,ktTeal,line width=2.4pt] (b)--(h);
  \node[text=ktInk,font=\footnotesize,fill=white,inner sep=2pt] at (0,1.00)
      {$\phi\colon S^1\times D^2\hookrightarrow S^3$};
  \node[align=center,text=ktInk,font=\small] at (0,-1.18)
    {węzeł przyczepienia $K\subset S^3$ i ramowanie $n\in\mathbb Z$};
  \node[rounded corners=4pt,fill=ktGold!16,inner sep=6pt,
        text=ktInk,font=\small] at (0,-2.0)
    {$\partial W_n(K)\cong S^3_n(K)$};
\end{tikzpicture}
\end{center}

\section{Forma przecięcia a macierz splecenia}
\label{sec:kt-intersection}

\begin{twierdzenie}[Forma przecięcia z diagramu uchwytowego]
\label{thm:kt-handle-intersection}
Niech $L=K_1\cup\dots\cup K_r\subset S^3$ będzie
zorientowanym splotem z całkowitymi ramowaniami
$n_1,\dots,n_r$. Niech $W_L$ będzie gładką, zwartą,
orientowaną czterorozmaitością otrzymaną z $B^4$
przez dołączenie do niego $r$ uchwytów indeksu dwa
wzdłuż $(L;n_1,\dots,n_r)$. Wtedy:
\begin{enumerate}[label=(\roman*)]
\item $\pi_1(W_L)=1$;
\item $H_2(W_L;\mathbb Z)\cong\mathbb Z^r$;
\item w bazie wyznaczonej przez rdzenie uchwytów po
uzupełnieniu ich powierzchniami Seiferta forma przecięcia
$Q_{W_L}$ jest reprezentowana przez macierz
$Q_L$ o przekątnej $(n_1,\dots,n_r)$ i elementach
pozadiagonalnych $\operatorname{lk}(K_i,K_j)$.
\end{enumerate}
\end{twierdzenie}

\begin{proof}
Zacznijmy od grupy fundamentalnej. Pogrubiając
sklejane części kołnierzami, stosujemy twierdzenie
van Kampena do otwartego pokrycia, w którym jedna
część jest równoważna homotopijnie poprzedniej
rozmaitości, druga kontraktowalnemu uchwytowi,
a przecięcie obszarowi przyczepienia
$S^1\times D^2$. Grupa tego przecięcia jest
generowana przez klasę krzywej przyczepienia.
Twierdzenie van Kampena daje więc iloraz grupy
poprzedniej rozmaitości przez normalne domknięcie
tej klasy. Ponieważ $\pi_1(B^4)=1$,
po wszystkich dołączeniach nadal mamy $\pi_1(W_L)=1$.

Następnie skorzystajmy z względnego kompleksu komórkowego
rozkładu na uchwyty pary $(W_L,B^4)$.
W tej parze występuje dokładnie $r$ komórek wymiaru dwa,
nie ma natomiast względnych komórek wymiarów jeden i trzy.
Istotnie, para uchwytu i obszaru przyczepienia
$(D^2\times D^2,\partial D^2\times D^2)$ jest
równoważna homotopijnie parze $(D^2,\partial D^2)$
przez ściąganie drugiego czynnika do zera.
Kompleks uchwytowy ma zatem
$C_2(W_L,B^4;\mathbb Z)\cong\mathbb Z^r$,
$C_1(W_L,B^4;\mathbb Z)=C_3(W_L,B^4;\mathbb Z)=0$,
a wszystkie różniczki są zerowe. Zatem
$H_2(W_L,B^4;\mathbb Z)\cong\mathbb Z^r$.
W odpowiednim fragmencie ciągu długiego pary mamy
\[
H_2(B^4)\longrightarrow H_2(W_L)
\longrightarrow H_2(W_L,B^4)
\longrightarrow H_1(B^4).
\]
Skrajne grupy są zerowe, więc środkowe odwzorowanie
jest izomorfizmem. Otrzymujemy
$H_2(W_L;\mathbb Z)\cong\mathbb Z^r$.

Aby obliczyć przecięcia, dla każdego $K_i$ wybieramy
powierzchnię Seiferta $F_i\subset S^3$. Jej wnętrze
przesuwamy nieznacznie do wnętrza $B^4$, a następnie
dołączamy rdzeń $D^2\times\{0\}$ odpowiedniego uchwytu.
Po wygładzeniu narożnika otrzymujemy zamkniętą
zorientowaną powierzchnię $\widehat F_i$.
Jej klasa przechodzi na $i$-ty generator
$H_2(W_L,B^4;\mathbb Z)$, ponieważ część w $B^4$
znika w homologii względnej, a rdzeń reprezentuje
generator względnej homologii uchwytu.
Klasy $e_i=[\widehat F_i]$ tworzą więc bazę.

Przypomnijmy rachunek przecięć w kuli.
Dla rozłącznych zorientowanych krzywych $J_1,J_2$
w $S^3$, gdy właściwie zanurzone zorientowane
powierzchnie $G_1,G_2\subset B^4$ mają te krzywe
za brzegi, ich algebraiczna liczba przecięć
wynosi $\operatorname{lk}(J_1,J_2)$.
Aby to zobaczyć, wybieramy powierzchnię Seiferta
$J_1$ w $S^3$ poprzeczną do $J_2$. Wpychamy ją
na płytki poziom kołnierza, zachowując brzeg,
a $G_2$ ustawiamy w postaci iloczynowej przy brzegu.
Przecięcia zachodzą wówczas dokładnie tam, gdzie
powierzchnia Seiferta przecina $J_2$; orientacja
brzegowa $S^3$ i orientacja kołnierza dają te same
znaki. Liczba nie zależy od wyboru powierzchni
w kuli: po uzgodnieniu kołnierzy dwóch
reprezentantów o tym samym brzegu ich różnica
jest 2-cyklem w $\operatorname{int}B^4$.
Ogranicza tam 3-łańcuch, gdyż wnętrze kuli
jest kontraktowalne. Wzór na brzeg przecięcia
z drugą powierzchnią pokazuje, że zmiana
algebraicznej liczby przecięć jest zerowa
(3-łańcuch nie spotyka jej brzegu w $S^3$).

Dla $i\ne j$ rdzenie różnych uchwytów są rozłączne,
więc wszystkie przecięcia $\widehat F_i$
i $\widehat F_j$ liczymy w $B^4$.
Z powyższego rachunku wynika
$e_i\cdot e_j=\operatorname{lk}(K_i,K_j)$.
Dla samoprzecięcia zastępujemy rdzeń $i$-tego
uchwytu przez rozłączny równoległy dysk
$D^2\times\{z\}$ dla małego $z\ne0$.
Jego brzeg po przyczepieniu jest równoległą
kopią $K_i'$ o klasie $n_i[\mu_i]+[\lambda_i]$.
Domykamy ten dysk powierzchnią Seiferta $K_i'$
wpchniętą do kuli. Otrzymana powierzchnia
reprezentuje nadal $e_i$, gdyż ma tę samą klasę
w $H_2(W_L,B^4;\mathbb Z)$.
Rdzenie są rozłączne, a rachunek w kuli daje
$e_i\cdot e_i=\operatorname{lk}(K_i,K_i')=n_i$
z definicji długości preferowanej i meridianu.
Zatem macierz przecięć jest równa $Q_L$.
Por. R.~E. Gompf, A.~I. Stipsicz,
\emph{4-Manifolds and Kirby Calculus}, AMS, 1999,
stwierdzenie 4.5.1.
\end{proof}

\begin{wniosek}
Jeżeli $W_L$ jest jak wyżej, to $\partial W_L$ jest
homologiczną sferą trójwymiarową, tzn.
$H_*(\partial W_L;\mathbb Z)\cong H_*(S^3;\mathbb Z)$,
wtedy i tylko wtedy, gdy $\det Q_L=\pm1$.
\end{wniosek}

\begin{proof}
W szczególności rozpiszmy zależność, której używamy.
Dla pary $(W_L,Y)$, gdzie $Y=\partial W_L$, mamy
fragment ciągu dokładnego
\[
H_2(W_L;\mathbb Z)\xrightarrow{j}
H_2(W_L,Y;\mathbb Z)\longrightarrow
H_1(Y;\mathbb Z)\longrightarrow H_1(W_L;\mathbb Z)=0.
\]
Dualność Poincarégo--Lefschetza identyfikuje
$H_2(W_L,Y;\mathbb Z)$ z $H^2(W_L;\mathbb Z)$,
a twierdzenie o uniwersalnych współczynnikach,
wobec $H_1(W_L;\mathbb Z)=0$, identyfikuje tę grupę
z $\operatorname{Hom}(H_2(W_L;\mathbb Z),\mathbb Z)$.
W bazie $e_i$ i bazie do niej dualnej mapa $j$
ma z definicji formy przecięcia macierz $Q_L$:
funkcjonał odpowiadający $j(e_i)$ przyjmuje
na $e_j$ wartość $e_i\cdot e_j$.
Jest to drugi dowód prezentacji homologii brzegu
z sekcji \ref{sec:kt-chirurgia-link}.
Zatem $H_1(Y;\mathbb Z)\cong\operatorname{coker}Q_L$.
Grupa ta znika dokładnie wtedy, gdy $Q_L$ jest
unimodularna (postać normalna Smitha).
Brzeg $Y$ jest spójny, gdyż powstaje przez
wypełnianie torusów brzegowych spójnego
zewnętrza splotu. Dla zamkniętej spójnej orientowalnej
trójrozmaitości dualność Poincarégo, uniwersalne twierdzenie
o współczynnikach i orientowalność pokazują, że z
$H_1=0$ wynika $H_2=0$, a $H_0=H_3=\mathbb Z$.
\end{proof}

\begin{przyklad}[Trzy klasyczne formy]
Dla pojedynczej składowej trywialnej o ramowaniu $+1$
mamy macierz przecięcia $Q=[1]$. Jej wyznacznik jest równy $1$,
więc brzeg jest homologiczną sferą trójwymiarową;
jest to nawet $S^3$, jak wykazano w sekcji
\ref{sec:kt-dehn} dla chirurgii $\pm1$ na węźle
trywialnym. Identyfikację wnętrza po domknięciu
uzasadniamy osobno, a nie przez samą macierz.

W $\mathbb{CP}^2$ prosta rzutowa $\mathbb{CP}^1$
ma normalną wiązkę zespoloną $\mathcal O(1)$.
Jej liczbę Eulera można odczytać również
z przecięcia dwóch różnych prostych rzutowych:
spotykają się poprzecznie w jednym punkcie
o dodatnim znaku. Otoczenie dyskowe prostej
jest więc wiązką dysków nad $S^2$ o liczbie
Eulera $+1$. Dzieląc bazę na dwie półsfery
dyskowe, przedstawiamy tę wiązkę jako kulę
$D^2\times D^2$ i jeden uchwyt indeksu dwa:
na równiku funkcja przejścia wiązki ma stopień
$+1$, co daje ramowanie $+1$ węzła trywialnego.
Dopełnienie odpowiednio wybranego otoczenia
prostej jest kulą: w współrzędnych jednorodnych
$[z_0:z_1:z_2]$ prosta $z_2=0$ ma otoczenie
$|z_2|^2\le\varepsilon^2(|z_0|^2+|z_1|^2)$.
W mapie $z_2\ne0$, o współrzędnych
$(z_0/z_2,z_1/z_2)$, dopełnienie wnętrza
tego otoczenia jest standardową domkniętą kulą.
Po przyklejeniu tej kuli jako uchwytu indeksu
cztery otrzymujemy zatem $\mathbb{CP}^2$.
Odwrócenie orientacji daje odpowiednio
ramowanie $-1$ i $\overline{\mathbb{CP}}^{\,2}$
z formą przecięcia $[-1]$.

Dla standardowego splotu Hopfa, zorientowanego tak,
by liczba splecenia wynosiła $+1$, z dwoma ramowaniami zero mamy
\[
Q=\begin{pmatrix}0&1\\1&0\end{pmatrix}.
\]
Macierz ta ma wyznacznik $-1$, a więc ponownie brzeg jest
homologiczną sferą trójwymiarową. Jej wartości własne wynoszą
$1$ i $-1$, zatem forma jest hiperboliczna i ma sygnaturę zero.
Brzeg jest w istocie $S^3$, zgodnie z przykładem
chirurgii na splocie Hopfa w sekcji
\ref{sec:kt-chirurgia-link}.
Rozkład uchwytowy $S^2\times S^2$ można otrzymać
bezpośrednio z rozkładów obu sfer na dysk indeksu
zero i dysk indeksu dwa. Iloczyn dwóch dysków
indeksu zero jest kulą $D^2\times D^2$.
Dwa iloczyny dysku indeksu dwa z dyskiem indeksu
zero stanowią uchwyty indeksu dwa.
Ich krzywe przyczepienia na brzegu początkowej
kuli to $S^1\times\{0\}$ i $\{0\}\times S^1$,
czyli splot Hopfa. Ramowania pochodzą od
struktury iloczynowej i są zerowe: normalne
wiązki dwóch sfer czynnikowych są trywialne,
więc każda ma samoprzecięcie zero. Pozostaje
iloczyn obu dysków indeksu dwa, czyli jeden
uchwyt indeksu cztery. Zatem w tym przykładzie
domknięcie przez jeden uchwyt indeksu cztery
daje $S^2\times S^2$, bez uchwytów indeksu trzy.

Ten przykład pokazuje, jak bezpośrednio z diagramu odczytać zarówno
homologię brzegu, bo $H_1(\partial W;\mathbb Z)\cong\operatorname{coker}Q$,
jak i informację o formie przecięcia we wnętrzu. Dla splotu Hopfa
z ramowaniami $n_1,n_2$ otrzymalibyśmy ogólniej macierz
$\begin{psmallmatrix}n_1&1\\1&n_2\end{psmallmatrix}$.
Jej wyznacznik wynosi $d=n_1n_2-1$.
Ponieważ w macierzy występuje wpis $1$,
pierwszy dzielnik determinantalny wynosi $1$.
Postać Smitha ma przekątną $(1,|d|)$, także
dla $d=0$. Zatem pierwsza homologia brzegu
jest izomorficzna z $\mathbb Z/(d)$ (dla
$d=0$ jest to $\mathbb Z$).
Jeżeli $d<0$, wartości własne mają przeciwne
znaki i sygnatura wynosi zero. Gdy $d>0$,
mają ten sam znak, wyznaczony przez ślad
$n_1+n_2$, więc sygnatura wynosi
$2\operatorname{sgn}(n_1+n_2)$. Dla $d=0$
jedna wartość własna jest zerowa, a druga
równa niezerowemu śladowi; sygnatura wynosi
wtedy $\operatorname{sgn}(n_1+n_2)$.
\end{przyklad}

\section{Przesuwanie uchwytów i operacja blow-up}
\label{sec:kt-slides}

\begin{stwierdzenie}[Przesunięcie uchwytu]
\label{st:kt-handle-slide}
Niech $W_L$ będzie czterorozmaitością zbudowaną z kuli
$B^4$ przez dołączenie uchwytów indeksu dwa opisanych
przez ramowany splot $L$ z macierzą przecięcia $Q_L$.
Przesunięcie jednego uchwytu po innym realizuje zmianę
bazy w $H_2(W_L;\mathbb Z)$ postaci
$e_i\mapsto e_i\pm e_j$, pozostawiając pozostałe
wektory bez zmian. Odpowiada to kongruencji całkowitej
$Q_L\mapsto P^TQ_LP$ dla pewnej macierzy
$P\in GL_r(\mathbb Z)$.
\end{stwierdzenie}

\begin{proof}
Korzystamy z geometrycznego twierdzenia o
przesuwaniu uchwytów tego samego indeksu:
zob. R.~E. Gompf, A.~I. Stipsicz,
\emph{4-Manifolds and Kirby Calculus}, AMS, 1999,
opis przesuwania uchwytów indeksu dwa.
Twierdzenie to zapewnia dyfeomorfizm z zachowaniem
orientacji między rozmaitościami przed i po
przesunięciu, indukujący dyfeomorfizm brzegów
(nie twierdzimy, że jest identycznością na brzegu).

Pod tym utożsamieniem nowy względny rdzeń
$i$-tego uchwytu reprezentuje sumę starego
rdzenia i równoległego rdzenia $j$-tego,
połączonych pasmem w kołnierzu. Znak drugiego
składnika $\varepsilon\in\{1,-1\}$ zależy
od orientacji kopii równoległej. Po domknięciu
w $B^4$ otrzymujemy $e_i'=e_i+\varepsilon e_j$,
co wynika z izomorfizmu homologii bezwzględnej
i względnej wykazanego w twierdzeniu
\ref{thm:kt-handle-intersection}.
Przy konwencji kolumnowej macierz przejścia
to $P=I+\varepsilon E_{ji}$, gdzie $i\ne j$,
a $E_{ji}$ ma jedynkę w pozycji $(j,i)$
i zera poza nią. Ponieważ $E_{ji}^2=0$,
$P^{-1}=I-\varepsilon E_{ji}$, a $\det P=1$.
Biliniowość daje
$(Q')_{ab}=(e_a')\cdot(e_b')=(P^TQ_LP)_{ab}$.
W szczególności sygnatura, rząd i kokernel
formy pozostają niezmienione. Zgodność form
splecenia na brzegu sprawdzono już w sekcji
\ref{sec:kt-linking-pairing}.
\end{proof}

\begin{przyklad}[Zmiana wszystkich współczynników przy przesunięciu]
Niech dwie składowe ramowanego splotu mają ramowania
$n_i,n_j$ i liczbę splecenia $\ell$. W bazie klas
rdzeni uchwytów macierz przecięcia wynosi
$Q=\begin{psmallmatrix}n_i&\ell\\\ell&n_j\end{psmallmatrix}$.
Przesunięcie $i$-tego uchwytu po $j$-tym odpowiada
zamianie $e_i'=e_i+e_j$, $e_j'=e_j$, a więc
\[
Q'=\begin{pmatrix}n_i+2\ell+n_j&\ell+n_j\\
\ell+n_j&n_j\end{pmatrix}.
\]
Nowe ramowanie $i$-tej składowej nie jest zatem zwykłą
sumą $n_i+n_j$: pojawia się dodatkowy składnik
$2\operatorname{lk}(K_i,K_j)$. Rachunek wynika z
równości $(e_i+e_j)^2=e_i^2+2(e_i\cdot e_j)+e_j^2$.
\end{przyklad}

\begin{stwierdzenie}[Blow-up]
Niech $W$ będzie zwartą, spójną, gładką,
orientowaną czterorozmaitością, ewentualnie z brzegiem.
Formę przecięcia rozpatrujemy na grupie
\[
H_2(W;\mathbb Z)/\operatorname{Tor}H_2(W;\mathbb Z),
\]
gdzie $\operatorname{Tor}$ oznacza podgrupę torsyjną.
Po wykonaniu we wnętrzu $W$ operacji blow-up otrzymujemy
$W\#\overline{\mathbb{CP}}^{\,2}$; po zmianie znaku
operacji odpowiednio $W\#\mathbb{CP}^2$.
Forma przecięcia zmienia się przez sumę prostą
$Q_W\oplus[-1]$ albo $Q_W\oplus[1]$,
niezależnie od tego, czy forma $Q_W$ jest
niezdegenerowana. Brzeg z orientacją pozostaje
niezmieniony. W szczególności jest to operacja
zachowująca trójrozmaitość brzegową, ale nie
typ dyfeomorficzny czterorozmaitości: zwiększa
jej drugą liczbę Bettiego o jeden.
\end{stwierdzenie}

\begin{proof}
Usuwamy z wnętrza $W$ standardową kulę $B^4$ i wklejamy
$\mathbb{CP}^2\setminus\operatorname{int}B^4$ lub
$\overline{\mathbb{CP}}^{\,2}\setminus\operatorname{int}B^4$.
Sklejenie odwraca orientacje obu sfer brzegowych.
Definicja sumy spójnej daje wskazany opis,
a operacja we wnętrzu nie zmienia pierwotnego brzegu.
Z wycinania i homologii pary $(B^4,S^3)$ wynika,
że usunięcie wnętrza kuli nie zmienia $H_2$
żadnego składnika: względne grupy w stopniach
dwa i trzy są zerowe. Następnie stosujemy
Mayera--Vietorisa (z pogrubieniem kołnierzy):
\[
0=H_2(S^3)\to H_2(W)\oplus H_2(\mathbb{CP}^2)
\to H_2(W\#\mathbb{CP}^2)\to H_1(S^3)=0.
\]
Dla odwróconej orientacji rozumujemy tak samo.
Po przejściu do ilorazów przez torsję otrzymujemy
sumę prostą grup wolnych. Reprezentanci klas
z obu składników mogą być wybrani z dala
od kuli i sfery sklejania; ich wzajemne
przecięcia są zerowe, a przecięcia w obrębie
każdego składnika nie zmieniają się. Formy
na $\mathbb{CP}^2$ i $\overline{\mathbb{CP}}^{\,2}$
to odpowiednio $[1]$ i $[-1]$, jak w przykładzie
poprzedniej sekcji. Wynikowa forma jest zatem
sumą prostą, a druga liczba Bettiego wzrasta
o jeden. Z twierdzenia van Kampena wynika
ponadto, że grupa fundamentalna się nie zmienia:
usunięcie kuli nie zmienia $\pi_1$, a drugi
składnik i sfera sklejania są jednospójne.
\end{proof}

\begin{uwaga}[Granica rachunku macierzowego]
Dwie macierze przecięcia całkowicie izometryczne nie muszą
pochodzić od dyfeomorficznych czterorozmaitości.
Rachunek na macierzach zapomina dodatkowe dane zanurzeń,
gładkości i uchwytów wyższych indeksów. To jeden z powodów,
dla których wymiar cztery jest szczególny: klasyfikacja
samych form przecięcia jest znacznie łatwiejsza niż
klasyfikacja gładkich czterorozmaitości.
\end{uwaga}

\section{Węzły jako brzegi powierzchni w \texorpdfstring{$B^4$}{B4}}
\label{sec:kt-slice}

\begin{definicja}[Rodzaj czterowymiarowy]
Niech $K\subset S^3=\partial B^4$ będzie zorientowanym węzłem.
\emph{Gładkim rodzajem czterowymiarowym} węzła nazywamy
\[
\begin{aligned}
g_4^{\mathrm{sm}}(K):=\min\big\{g(F):\;&F\subset B^4
\text{ jest spójną, zwartą,}\\
&\text{orientowaną, gładko i właściwie zanurzoną powierzchnią},\\
&\partial F=K\big\}.
\end{aligned}
\]
Analogicznie \emph{topologiczny rodzaj czterowymiarowy}
$g_4^{\mathrm{top}}(K)$ definiujemy, dopuszczając powierzchnie
topologicznie lokalnie płasko i właściwie zanurzone.
Węzeł jest \emph{gładko slice}, gdy $g_4^{\mathrm{sm}}(K)=0$,
a \emph{topologicznie slice}, gdy $g_4^{\mathrm{top}}(K)=0$.
\end{definicja}

\begin{stwierdzenie}
Dla każdego węzła $K\subset S^3$ zachodzi
$0\leq g_4^{\mathrm{top}}(K)\leq g_4^{\mathrm{sm}}(K)\leq g(K)$.
\end{stwierdzenie}

\begin{proof}
Każda powierzchnia gładko zanurzona jest lokalnie płaska,
więc minimum topologiczne nie przekracza gładkiego.
Każdą powierzchnię Seiferta w $S^3$ można przesunąć
nieznacznie w głąb kołnierza $S^3\times[0,\varepsilon)$
we wnętrzu $B^4$, pozostawiając jej brzeg na $S^3$.
Dostajemy właściwie zanurzoną powierzchnię tego samego
rodzaju, stąd $g_4^{\mathrm{sm}}(K)\leq g(K)$.
Dokładniej, w kołnierzu obraz $F$ zastępujemy
wykresem gładkiej funkcji $u:F\to[0,\varepsilon)$,
dodatniej we wnętrzu, zerowej na brzegu i mającej
niezerową pochodną w kierunku do wnętrza przy brzegu.
Wykres jest zanurzeniem, ponieważ rzut na $F$
jest jego odwrotnością, a warunek na pochodną
zapewnia poprzeczność do $S^3$ wzdłuż $K$.
Oba zbiory dopuszczalnych rodzajów są zatem
niepustymi podzbiorami $\mathbb N_0$, więc mają
elementy najmniejsze. To uzasadnia też istnienie
minimów w definicji.
\end{proof}

\begin{definicja}[Konkordancja węzłów]
Niech $K_0,K_1\subset S^3$ będą zorientowanymi węzłami.
\emph{Gładką konkordancją} między $K_0$ i $K_1$ nazywamy
gładkie, właściwe zanurzenie
$A\colon S^1\times[0,1]\hookrightarrow S^3\times[0,1]$
takie, że $A(S^1\times\{i\})=K_i\times\{i\}$ dla $i=0,1$
oraz orientacja pierścienia indukuje na górnym
brzegu orientację $K_1$, a na dolnym orientację
przeciwną do orientacji $K_0$.
Jest to szczególny rodzaj kobordyzmu: powierzchnia
łącząca węzły jest pierścieniem, a nie dowolną
powierzchnią. Zakładamy produktową postać przy
brzegu (można ją uzyskać przez wyprostowanie
kołnierza). Analogicznie definiujemy
\emph{konkordancję topologiczną} przez lokalnie
płaskie właściwe zanurzenie pierścienia,
z lokalnie płaskim kołnierzem przy brzegu.
W dalszej części symbol $-K$ oznacza odbicie
lustrzane węzła z odwróceniem jego orientacji,
nie samo odwrócenie orientacji krzywej.
\end{definicja}

\begin{stwierdzenie}[Konkordancja a dysk w kuli]
\label{st:kt-concordance-slice}
Dwa zorientowane węzły $K_0,K_1\subset S^3$ są gładko
konkordantne wtedy i tylko wtedy, gdy
$K_0\#(-K_1)$ jest gładko slice, gdzie $-K_1$
oznacza odbicie lustrzane $K_1$ z odwróconą orientacją.
Analogiczna równoważność zachodzi w kategorii topologicznej
lokalnie płaskiej.
\end{stwierdzenie}

\begin{proof}
Najpierw w kategorii gładkiej rozpatrzmy trzy
konstrukcje.

\emph{Dysk i konkordancja do węzła trywialnego.}
Jeżeli $D\subset B^4$ jest właściwie zanurzonym
dyskiem, wybieramy małą kulę wokół punktu
wewnętrznego $D$, w lokalnych współrzędnych pary
$(B^4,D)$. Przecina ona $D$ w standardowym dysku.
Kulę tę można wybrać jako obraz małej standardowej
kuli przez dyfeomorfizm o nośniku w otoczeniu
punktu. Uzupełnienie jej wnętrza w $B^4$ jest
zatem dyfeomorficzne z $S^3\times[0,1]$.
Pozostała część $D$ jest pierścieniem łączącym
$K$ z węzłem trywialnym. Odwrotnie, do końca
konkordancji przy węźle trywialnym doklejamy
standardową kulę z jej standardowym dyskiem.
Po wygładzeniu narożników otrzymujemy dysk
w $B^4$. Węzeł ogranicza więc dysk wtedy
i tylko wtedy, gdy jest konkordantny z trywialnym.

\emph{Rozcięcie pierścienia.}
W pierścieniu konkordancji wybieramy właściwie
zanurzony łuk łączący oba brzegi, produktowy
przy końcach. Po izotopii otoczenia możemy
przyjąć, że łuk jest $\{p\}\times[0,1]$.
Uzasadnienie: dwa łuki o tych samych końcach
są homotopijne względem końców w jednospójnym
$S^3\times[0,1]$. Po ustaleniu małych
końcowych odcinków ogólne położenie homotopii
w wymiarze cztery pozwala wybrać ją jako
rodzinę zanurzeń łuku bez samoprzecięć:
parametry dwóch punktów łuku i czasu mają
łącznie wymiar trzy, mniejszy od kowymiaru
przekątnej w iloczynie dwóch kopii
czterorozmaitości. Warunek niezerowej
pochodnej wzdłuż łuku uzyskuje się również
przez transwersalność. Twierdzenie o
rozszerzaniu izotopii daje żądaną izotopię
otoczenia. Punkty na obu końcach możemy
uprzednio utożsamić przez izotopie sfer.
Wybieramy małe produktowe otoczenie łuku
$D^3\times[0,1]$, w którym powierzchnia
przecina je w pasie $[-1,1]\times[0,1]$.
Taki model wynika z otoczenia tubularnego
łuku w parze (normalne wiązki nad przedziałem
są trywialne). Po usunięciu wnętrza otoczenia
otrzymujemy $(S^3\setminus\operatorname{int}D^3)
\times[0,1]\cong D^3\times[0,1]\cong B^4$
(po wygładzeniu narożników).
Pierścień po usunięciu pasa jest dyskiem.
Na brzegu kuli widzimy pozostałe łuki dwóch
węzłów połączone wzdłuż bocznego brzegu.
Dolna sfera ma przeciwną orientację otoczenia
względem górnej, a dolny brzeg pierścienia
ma odwróconą orientację krzywej.
Brzeg dysku jest zatem $K_1\#(-K_0)$.
Po odwróceniu kierunku konkordancji
otrzymujemy $K_0\#(-K_1)$.

\emph{Suma spójna konkordancji.}
Dla dwóch konkordancji rozcinamy ich pierścienie
oraz otoczenia wzdłuż takich łuków, jak wyżej.
Sklejamy boczne części $\partial D^3\times[0,1]$
tak, aby sklejały się dwa brzegi pasów z
odpowiednimi orientacjami. Na poziomie otoczenia
jest to iloczyn sklejenia dwóch kul $D^3$
wzdłuż ich sfer brzegowych z przedziałem,
czyli ponownie $S^3\times[0,1]$.
Powierzchnia jest pierścieniem, ponieważ
sklejamy dwa dyski wzdłuż dwóch rozłącznych
łuków na ich brzegach, w sposób orientowalny.
Na poziomach końcowych dostajemy sumy spójne
odpowiednich węzłów. Suma spójna jest więc
dobrze określona na klasach konkordancji.

Jeżeli $K_0\#(-K_1)$ ogranicza dysk, jest
konkordantny z węzłem trywialnym. Z rozcięcia
produktowej konkordancji $K_1\times[0,1]$
wynika też, że $K_1\#(-K_1)$ ogranicza dysk.
Biorąc sumę spójną pierwszej konkordancji
z produktową konkordancją $K_1$ i korzystając
z łączności oraz przemienności sumy spójnej,
otrzymujemy konkordancję $K_0$ z $K_1$.
Nie zakładamy tu istnienia grupy konkordancji.

W kategorii topologicznej korzystamy z klasycznej
lokalnie płaskiej wersji tych konstrukcji,
w szczególności prostowania łuku w parze
i otoczeń produktowych; nie wyprowadzamy ich
z gładkiej transwersalności. Zob.
P.~Teichner, \emph{Slice Knots: Knot Theory in
the 4th Dimension}, notatki R.~Robertsa,
sekcja o grupie konkordancji węzłów.
Konstrukcje rozcięcia, sklejenia i argument
algebraiczny pozostają takie same, co daje
analogiczną równoważność.
\end{proof}

\begin{definicja}[Grupa konkordancji]
Niech $\mathcal C^{\mathrm{sm}}$ oznacza zbiór klas gładkiej
konkordancji zorientowanych węzłów w $S^3$.
Operacją jest suma spójna $[K_1]+[K_2]:=[K_1\#K_2]$.
Elementem neutralnym jest klasa węzła trywialnego,
a elementem przeciwnym do $[K]$ jest $[-K]$.
Tak zdefiniowana grupa abelowa nazywa się
\emph{gładką grupą konkordancji węzłów}.
\end{definicja}

\begin{proof}[Uzasadnienie struktury grupy]
Konkordancja jest relacją równoważności:
zwrotność daje pierścień produktowy, symetrię
otrzymujemy przez odwrócenie parametru przedziału
i orientacji pierścienia, a przechodniość przez
sklejenie wzdłuż wspólnego brzegu w produktowych
kołnierzach i przeparametryzowanie przedziału.
Dobra określoność sumy na klasach wynika z
konstrukcji sumy konkordancji w dowodzie
stwierdzenia \ref{st:kt-concordance-slice}.
Łączność, przemienność i element neutralny
wynikają z odpowiednich własności sumy spójnej
z dokładnością do izotopii. Rozcięcie produktowej
konkordancji $K\times[0,1]$ daje dysk o brzegu
$K\#(-K)$, więc na mocy tego samego stwierdzenia
$[K]+[-K]=[U]$, gdzie $U$ jest węzłem trywialnym.
To kończy dowód wszystkich aksjomatów grupy.
Analogicznie otrzymujemy grupę topologiczną
$\mathcal C^{\mathrm{top}}$. Każda gładka
konkordancja jest lokalnie płaska, więc
przyporządkowanie tej samej klasy w kategorii
topologicznej definiuje homomorfizm
$\mathcal C^{\mathrm{sm}}\to\mathcal C^{\mathrm{top}}$.
\end{proof}

\section{Węzły ribbon i nierówności rodzaju}
\label{sec:kt-ribbon}

\begin{definicja}[Węzeł ribbon]
Węzeł $K\subset S^3$ jest \emph{ribbon}, jeżeli ogranicza
immersję dysku w $S^3$, której wszystkie osobliwości
są poprzecznymi łukami podwójnymi typu wstęgowego:
przeciwobraz każdego takiego łuku składa się z jednego
łuku właściwie zanurzonego w dysku i jednego łuku leżącego
we wnętrzu dysku. Równoważnie $K$ ogranicza w $B^4$
gładki dysk, dla którego funkcja radialna jest funkcją
Morse'a bez lokalnych maksimów we wnętrzu dysku.
\end{definicja}

\begin{stwierdzenie}
Każdy węzeł ribbon jest gładko slice.
\end{stwierdzenie}

\begin{proof}
W zanurzonym dysku typu ribbon przesuwamy lokalnie jedną z gałęzi
przecinających się w stronę wnętrza $B^4$ zgodnie
z wysokością radialną. Ponieważ osobliwości mają typ wstęgowy,
przesunięcie eliminuje wszystkie łuki podwójne bez tworzenia
nowych, a otrzymany dysk jest gładko zanurzony w $B^4$.
Brzeg pozostaje równy $K$.
\end{proof}

\begin{uwaga}[Hipoteza slice--ribbon]
Odwrotność poprzedniego stwierdzenia, tzn. pytanie,
czy każdy gładko slice węzeł jest ribbon, to klasyczna
hipoteza Foxa. Jest to problem otwarty w klasycznej
teorii konkordancji w $S^3$.
Nie należy go mylić z implikacją topologicznie slice
$\Rightarrow$ gładko slice, która jest fałszywa.
\end{uwaga}

\begin{lemat}[Wymiar podprzestrzeni izotropowej a sygnatura]
\label{lem:kt-isotropic-signature}
Niech $V$ będzie rzeczywistą przestrzenią liniową
wymiaru $n$, a $S\colon V\times V\to\mathbb R$
niezdegenerowaną formą symetryczną o bezwładności
$(s_+,s_-)$, gdzie $n=s_++s_-$.
Jeżeli $L\subset V$ jest podprzestrzenią
\emph{całkowicie izotropową}, tzn.
$S(x,y)=0$ dla wszystkich $x,y\in L$,
to $\dim L\le\min(s_+,s_-)$.
W szczególności
\[
|\operatorname{sign}S|
\le n-2\dim L.
\]
\end{lemat}

\begin{proof}
Z prawa bezwładności Sylvestera wybierzmy rozkład
$V=V_+\oplus V_-$, na którym forma $S$
jest odpowiednio dodatnio i ujemnie określona,
a składniki są względem $S$ ortogonalne.
Rzutowanie $\pi_+\colon L\to V_+$
musi być injektywne: gdyby $x\in L$
miało $\pi_+(x)=0$, należałoby do $V_-$,
a z izotropowości $S(x,x)=0$ oraz
ujemnej określoności na $V_-$
wynikałoby $x=0$.
Analogicznie $\pi_-\colon L\to V_-$
jest injektywne. Stąd
$\dim L\le s_+$ i $\dim L\le s_-$.
Ostatnia nierówność wynika z równości
$|\operatorname{sign}S|=|s_+-s_-|
=s_++s_--2\min(s_+,s_-)$.
\end{proof}

\begin{twierdzenie}[Nierówność Murasugiego dla rodzaju czterowymiarowego]
Dla każdego zorientowanego węzła $K\subset S^3$
zachodzi $|\sigma(K)|\leq2g_4^{\mathrm{sm}}(K)$.
W szczególności, jeżeli $K$ jest gładko slice, to
$\sigma(K)=0$.
\end{twierdzenie}

\begin{proof}[Redukcja do podprzestrzeni izotropowej]
Niech $F$ będzie powierzchnią Seiferta węzła $K$ o
rodzaju $g$, a $G\subset B^4$ właściwie zanurzoną
zorientowaną powierzchnią o brzegu $K$ i rodzaju $h$.
Rozważmy symetryczną formę $S=V+V^T$ na
$H_1(F;\mathbb R)$, gdzie $V$ jest macierzą Seiferta.

Potrzebujemy klasycznego lematu Murasugiego:
\emph{przy powyższych założeniach istnieje rzeczywista
podprzestrzeń $L\subset H_1(F;\mathbb R)$
o wymiarze co najmniej $g-h$, na której forma $V$
znika identycznie}. Dowodzi się go przez połączenie
$F$ i $G$ w zamkniętą powierzchnię w $S^4$, wybranie
ograniczającej ją trójrozmaitości oraz zastosowanie
ciągu dokładnego pary i dualności Poincarégo--Lefschetza.
W przypadku $h=0$ daje on metabolizer opisany wcześniej.

Na $L$ znika również symetryzacja $S$. Ponieważ
$\det S=\pm\det(K)\ne0$, forma $S$ jest niezdegenerowana. Niech $s_+$ i $s_-$ oznaczają liczby dodatnich
i ujemnych wartości własnych $S$.
Ponieważ $S$ jest niezdegenerowana,
$s_++s_-=2g$. Z Lematu~\ref{lem:kt-isotropic-signature}
wynika
\[
|\operatorname{sign}S|
\le 2g-2\dim L\le2h.
\]
Minimalizując $h$ po wszystkich gładkich powierzchniach
w $B^4$, dostajemy $|\sigma(K)|\le2g_4^{\mathrm{sm}}(K)$.
\end{proof}

\begin{przyklad}[Oszacowanie rodzaju trójlistnika przez sygnaturę]
\label{ex:kt-trefoil-four-genus-signature}
Niech $T$ będzie prawoskrętnym trójlistnikiem.
Dla standardowej powierzchni Seiferta rodzaju jeden
można wybrać bazę, w której macierz symetryczna
formy Seiferta wynosi
\[
V+V^T=\begin{pmatrix}2&-1\\-1&2\end{pmatrix}.
\]
Jej wartości własne wynoszą $1$ i $3$,
zatem $|\sigma(T)|=2$.
Nierówność Murasugiego daje
$g_4^{\mathrm{sm}}(T)\ge1$.
Z drugiej strony powierzchnię Seiferta
rodzaju jeden można przesunąć do $B^4$,
więc $g_4^{\mathrm{sm}}(T)\le g(T)=1$.
Otrzymujemy
\[
g_4^{\mathrm{sm}}(T)=1.
\]
W szczególności trójlistnik nie jest
gładko slice. Ten rachunek jest
elementarnym przypadkiem formuły
dla rodzaju czterowymiarowego węzłów torusowych;
nie zastępuje jej ogólnego dowodu
dla dowolnych względnie pierwszych $p,q$.
\end{przyklad}

\begin{uwaga}
Ważne jest, że w twierdzeniu występuje rodzaj w $B^4$,
a nie tylko rodzaj powierzchni Seiferta w $S^3$.
Dowód wykorzystuje formy Seiferta i ograniczenie
na wymiar podprzestrzeni izotropowej pochodzącej
z powierzchni zanurzonej w $B^4$; jest istotnie głębszy
od prostego oszacowania $|\sigma(K)|\leq2g(K)$.
\end{uwaga}

\begin{twierdzenie}[Fox--Milnor]
Jeśli $K\subset S^3$ jest gładko slice, to istnieje
wielomian Laurenta $f(t)\in\mathbb Z[t,t^{-1}]$ taki, że
$\Delta_K(t)\doteq f(t)f(t^{-1})$,
gdzie $\doteq$ oznacza równość z dokładnością do
mnożenia przez $\pm t^k$.
\end{twierdzenie}

\begin{uwaga}
Warunek Foxa--Milnora nie jest wystarczający do gładkiej
slicowości. Ograniczenie to jest algebraiczne i nie wykrywa
wielu obstrukcji pochodzących z gładkiej topologii
czterowymiarowej. Zobaczymy niżej przykład o
$\Delta_K(t)=1$, który jest topologicznie slice,
ale nie jest gładko slice.
\end{uwaga}

\begin{definicja}[Metabolizer formy Seiferta]
Niech $F$ będzie powierzchnią Seiferta węzła $K\subset S^3$
o rodzaju $g$ i niech $V$ będzie macierzą jej formy Seiferta.
\emph{Metabolizerem} nazywamy podgrupę $L\subset H_1(F;\mathbb Z)$
która jest składnikiem prostym rzędu $g$ oraz na której
forma Seiferta znika identycznie, tj.
$\operatorname{lk}(x^+,y)=0$ dla wszystkich $x,y\in L$.
Węzeł nazywamy \emph{algebraicznie slice}, jeżeli jego
forma Seiferta dopuszcza metabolizer.
\end{definicja}

\begin{twierdzenie}[Slicowość gładka implikuje algebraiczną]
Niech $K\subset S^3$ będzie węzłem, który ogranicza gładko
zanurzony dysk $D\subset B^4$. Wtedy dla każdej powierzchni
Seiferta $F$ węzła $K$ forma Seiferta na $H_1(F;\mathbb Z)$
posiada metabolizer.
\end{twierdzenie}

\begin{proof}[Argument przez rozmaitość trójwymiarową]
Niech $F$ ma rodzaj $g$. Łącząc $F$ z dyskiem $D$
wzdłuż wspólnego brzegu, otrzymujemy zamkniętą,
zorientowaną powierzchnię $\Sigma$ rodzaju $g$ w
czterowymiarowej sferze powstałej po doklejeniu kuli
do $B^4$. Powierzchnia $\Sigma$ ogranicza
odpowiednią trójwymiarową rozmaitość Seiferta $M$;
w wersji dostosowanej do rozcięcia na dwie kule jej
brzeg jest sumą $F$ i dysku $D$ z przeciwnymi orientacjami.
Istnienie takiej rozmaitości $M$ jest klasycznym
lematem o hiperpłaskiej powierzchni Seiferta dla
zanurzenia kodymensji dwa.

Po zastosowaniu ciągu dokładnego pary $(M,\partial M)$
i dualności Poincarégo--Lefschetza otrzymujemy zasadę
\emph{połowa żyje, połowa ginie}: jądro
$H_1(\partial M;\mathbb Q)\to H_1(M;\mathbb Q)$
ma wymiar $g$ i jest lagranżowskie względem
formy przecięcia. Jej część całkowita, po wysyceniu
w $H_1(F;\mathbb Z)$, jest składnikiem prostym rzędu $g$.
Jeśli $x,y$ należą do tego składnika, to odpowiednie
wielokrotności krzywych ograniczają 2-łańcuchy w $M$.
Przesuwając jedną z nich w dodatnim kierunku normalnym
do $F$, można wybrać łańcuchy bez przecięć z drugą krzywą.
Stąd $\operatorname{lk}(x^+,y)=0$.
Otrzymany lagranżowski składnik jest metabolizerem.
\end{proof}

\begin{stwierdzenie}[Algebraiczna przyczyna faktoryzacji Foxa--Milnora]
Niech $K\subset S^3$ będzie algebraicznie slice,
a $V$ macierzą Seiferta rzędu $2g$.
Wtedy istnieje $f(t)\in\mathbb Z[t,t^{-1}]$ takie, że
$\det(tV-V^T)\doteq f(t)f(t^{-1})$.
\end{stwierdzenie}

\begin{proof}
Wybieramy całkowitą bazę $H_1(F;\mathbb Z)$,
której pierwsze $g$ elementów tworzy bazę metabolizera.
W tej bazie macierz Seiferta ma postać blokową
$V=\begin{psmallmatrix}0&A\\B&C\end{psmallmatrix}$,
gdzie bloki są kwadratowe rozmiaru $g$.
Stąd
\[
tV-V^T=\begin{pmatrix}
0&tA-B^T\\tB-A^T&tC-C^T
\end{pmatrix}.
\]
Rozwinięcie wyznacznika względem pierwszych $g$ wierszy
daje
$\det(tV-V^T)=(-1)^g\det(tA-B^T)\det(tB-A^T)$.
Niech $f(t):=\det(tA-B^T)$. Z transpozycji wynika
$\det(tB-A^T)=(-t)^g f(t^{-1})$.
Ostatecznie
$\Delta_K(t)=t^g f(t)f(t^{-1})$,
co dowodzi tezy z dokładnością do jednostki $t^g$.
\end{proof}

\begin{przyklad}[Jawna faktoryzacja Foxa--Milnora]
\label{ex:kt-fox-milnor-matrix}
Rozważmy całkowitą macierz
\[
V=\begin{pmatrix}0&2\\1&0\end{pmatrix}.
\]
Ponieważ
$V-V^T=\begin{psmallmatrix}0&1\\-1&0\end{psmallmatrix}$,
jej część antysymetryczna jest
unimodularną formą przecięcia.
Pierwszy wektor bazowy rozpina
metabolizer formy Seiferta, gdyż
$V_{11}=0$. Bezpośrednio obliczamy
\[
\begin{aligned}
\det(tV-V^T)
&=-\,(2t-1)(t-2)\\
&=t(2t-1)(2t^{-1}-1)
=-2t^2+5t-2.
\end{aligned}
\]
Wielomian ma więc wymaganą postać
$t\,f(t)f(t^{-1})$ z $f(t)=2t-1$.
W szczególności przy $t=1$ ma wartość $1$.
Jest to przykład \emph{rachunku algebraicznego}:
sama faktoryzacja nie wykazuje istnienia
gładkiego dysku typu slice dla węzła
o takiej macierzy.
\end{przyklad}

\begin{uwaga}
Klasyczna teoria algebraicznej konkordancji pokazuje,
że gładko slice węzeł jest algebraicznie slice.
Samo istnienie metabolizera nie gwarantuje jednak
zanurzonego dysku w $B^4$. Wyprowadzone twierdzenie
odtwarza zatem algebraiczną część przeszkody Foxa--Milnora,
ale nie jej odwrotność.
\end{uwaga}


\section{Twierdzenie Freedmana i rozdzielenie kategorii TOP/DIFF}
\label{sec:kt-freedman}

\begin{twierdzenie}[Freedman; trywialny wielomian Alexandera]
Niech $K\subset S^3$ będzie węzłem spełniającym
$\Delta_K(t)=1$ (w symetrycznej normalizacji).
Wówczas istnieje lokalnie płasko zanurzony dysk
$D^2\hookrightarrow B^4$ o brzegu $K$. Innymi słowy,
$g_4^{\mathrm{top}}(K)=0$.
\end{twierdzenie}

\begin{uwaga}
Dowód wykorzystuje topologiczną teorię zanurzania dysków,
uchwyty Cassona i czterowymiarową topologiczną chirurgię.
Nie otrzymujemy stąd gładko zanurzonego dysku. Wymiar
cztery jest szczególny właśnie dlatego, że topologicznie
dopuszczalne operacje na dyskach nie muszą mieć odpowiedników
gładkich.
\end{uwaga}

\begin{przyklad}[Podwojenie Whiteheada]
Dla dowolnego węzła $J\subset S^3$ jego nieskręcony
(zerowo ramowany) podwójny węzeł Whiteheada $\operatorname{Wh}(J)$
ma wielomian Alexandera $\Delta_{\operatorname{Wh}(J)}(t)=1$.
Z twierdzenia Freedmana jest więc topologicznie slice.
Dla dodatniego nieskręconego podwojenia prawoskrętnego
trójlistnika $T_{2,3}$ zachodzi jednak
$\tau(\operatorname{Wh}^{+}(T_{2,3}))=1$,
gdzie $\tau$ jest niezmiennikiem Floera opisanym w sekcji \ref{sec:kt-hfk}
rozdziału \ref{ch:kt-nowoczesne}. Z nierówności $|\tau(K)|\leq g_4^{\mathrm{sm}}(K)$
wynika, że ten węzeł nie jest gładko slice, choć jest
topologicznie slice.
\end{przyklad}

\begin{uwaga}[Znaczenie dla egzotycznych $\mathbb R^4$]
Rozbieżność między topologicznymi i gładkimi dyskami
w $B^4$ wiąże się z istnieniem niezgodnych struktur
gładkich na rozmaitościach homeomorficznych.
Szczegółowe konstrukcje egzotycznych $\mathbb R^4$
wymagają dodatkowo teorii uchwytów Cassona, twierdzeń
Donaldsona i innych narzędzi czterowymiarowych.
Sam fakt $g_4^{\mathrm{top}}\ne g_4^{\mathrm{sm}}$
jest ważnym, ale niepełnym opisem tej problematyki.
\end{uwaga}

\section{Węzły torusowe w \texorpdfstring{$B^4$}{B4} i problem Milnora}
\label{sec:kt-milnor}

\begin{twierdzenie}[Rozstrzygnięcie hipotezy Milnora]
Niech $p,q\geq2$ będą względnie pierwsze i niech
$T(p,q)\subset S^3$ będzie dodatnim węzłem torusowym.
Wtedy
$g_4^{\mathrm{sm}}(T(p,q))=(p-1)(q-1)/2$.
\end{twierdzenie}

\begin{uwaga}
Równość $g(T(p,q))=(p-1)(q-1)/2$ została udowodniona
wcześniej metodami klasycznymi. Znacznie trudniejsza
jest równość z rodzajem czterowymiarowym: a priori
powierzchnia w $B^4$ mogłaby mieć mniejszy rodzaj niż
powierzchnia w $S^3$. Pierwszy dowód dla całej rodziny
pochodzi z teorii cechowania Kronheimera--Mrowki.
Inne dowody wykorzystują teorię Seiberga--Wittena,
homologię Floera lub niezmiennik Rasmussena.
\end{uwaga}

\begin{przyklad}
Dla $T(2,3)$ mamy $g_4^{\mathrm{sm}}=1$;
dla $T(3,4)$ mamy $g_4^{\mathrm{sm}}=3$;
dla $T(3,5)$ mamy $g_4^{\mathrm{sm}}=4$.
Jest to rodzina, w której klasyczny rodzaj powierzchni
Seiferta i gładki rodzaj czterowymiarowy są sobie równe.
\end{przyklad}



% Iteracja 7: rozdzial autorski o plumbing, formach przeciecia i E8.
